% Restored from ../methods_hybrid_evidence_mil.tex and checked against
% ../../scripts/ss_experiment_suite.py: GatedAttentionPool,
% HybridEvidenceMILEncoder and SlideClassifier. This file is maintained by hand.
\subsection{Hybrid-Evidence MIL}

\paragraph{Task formulation and feature projection.}
Each WSI was represented by a sampled feature bag
$X=\{\mathbf{x}_i\}_{i=1}^{N}$ and its associated patch coordinates
$(a_i,b_i)$, where $N=512$ and $\mathbf{x}_i\in\mathbb{R}^{1536}$.
The slide label was $y=1$ for MALT and $y=0$ for SS.
No patch-level labels were used. A shared projection placed all features in a
$D=256$ dimensional hidden space:
\begin{equation}\label{eq:projection}
 \mathbf{h}_i = \operatorname{Dr}_{0.3}\!\left[
  \operatorname{GELU}\!\left(\operatorname{LN}
  (\mathbf{W}_{p}\mathbf{x}_i+\mathbf{b}_{p})\right)\right],
 \qquad \mathbf{h}_i\in\mathbb{R}^{D}.
\end{equation}
Here $\operatorname{LN}$ denotes layer normalization and $\operatorname{Dr}_{0.3}$
denotes dropout with probability 0.3 during training; dropout is disabled during
inference. All vectors below are column vectors, and $[\,\cdot\,;\,\cdot\,]$
denotes concatenation along the feature dimension.

\paragraph{Gated attention evidence.}
Gated pooling followed the attention mechanism of Ilse and colleagues
\citep{ref5}, with an additional value transformation implemented in the archived
model. For a sequence $Z=\{\mathbf{z}_j\}_{j=1}^{L}$ of $D$ dimensional vectors,
define the pooling operator $\mathcal{G}_{\theta}$ as
\begin{subequations}\label{eq:gated-pool}
\begin{align}
 \mathbf{t}_j &= \operatorname{Dr}_{0.3}
     \!\left[\operatorname{ReLU}(\mathbf{W}_{v}\mathbf{z}_j+\mathbf{b}_{v})\right],
     \label{eq:pool-value}\\
 \ell_j &= \mathbf{w}_{a}^{\mathsf T}
    \left[\tanh(\mathbf{V}\mathbf{t}_j+\mathbf{b}_{V})
     \odot\operatorname{sigmoid}(\mathbf{U}\mathbf{t}_j+\mathbf{b}_{U})\right]+b_a,
     \label{eq:pool-logit}\\
 \alpha_j &= \frac{\exp(\ell_j)}{\sum_{l=1}^{L}\exp(\ell_l)},
 \qquad \mathcal{G}_{\theta}(Z)=\sum_{j=1}^{L}\alpha_j\mathbf{t}_j.
     \label{eq:pool-sum}
\end{align}
\end{subequations}
The symbol $\odot$ denotes element-wise multiplication. All attention and value
matrices map between $D$ dimensional spaces. The attention evidence vector was
$\mathbf{e}_{\mathrm{att}}=\mathcal{G}_{\theta_A}(\{\mathbf{h}_i\}_{i=1}^{N})$.
Thus, the pooling weights multiply transformed values $\mathbf{t}_i$.
The regional branch used the same operator with independently learned parameters
$\theta_R$.

\paragraph{Distribution evidence.}
The center and heterogeneity of the sampled bag were summarized by two separate
vectors:
\begin{subequations}\label{eq:distribution}
\begin{align}
 \mathbf{e}_{\mu} &= \frac{1}{N}\sum_{i=1}^{N}\mathbf{h}_i,\\
 \mathbf{e}_{\sigma} &=
 \left[\frac{1}{N-1}\sum_{i=1}^{N}
       (\mathbf{h}_i-\mathbf{e}_{\mu})^{\odot 2}\right]^{\odot 1/2}.
\end{align}
\end{subequations}
Powers and square roots in Eq.~\eqref{eq:distribution} are element-wise. The
sample denominator $N-1$ matches the standard deviation used in the code.

\paragraph{Prototype-margin evidence.}
Two latent prototype banks indexed by $c\in\{0,1\}$ each contained
$K=8$ learnable prototypes
$\mathbf{p}_{c,r}\in\mathbb{R}^{D}$. Let a hat denote the implemented
$\ell_2$ normalization, with its numerical safeguard for zero norms.
The maximum cosine similarity within each bank and the bank-1-minus-bank-0 margin
were
\begin{subequations}\label{eq:prototype}
\begin{align}
 s_{i,c} &= \max_{1\leq r\leq K}
       \widehat{\mathbf{h}}_i^{\mathsf T}\widehat{\mathbf{p}}_{c,r},
       \qquad m_i=s_{i,1}-s_{i,0},\label{eq:margin}\\
 k &= \min(32,N),\qquad
 I_{+}=\operatorname{TopKIdx}(\{m_i\},k),\qquad
 I_{-}=\operatorname{TopKIdx}(\{-m_i\},k),\\
 \mathbf{e}_{+} &= \frac{1}{k}\sum_{i\in I_{+}}\mathbf{h}_i,
 \qquad \mathbf{e}_{-}=\frac{1}{k}\sum_{i\in I_{-}}\mathbf{h}_i.
\end{align}
\end{subequations}
$\operatorname{TopKIdx}$ returns indices of the $k$ largest entries. The high-
and low-margin vectors therefore average the instances with the largest and
smallest margins, respectively; the selected values need not all have the same
sign. Prototypes were optimized through differentiable similarity and margin
statistics in the slide-level classification objective. The hard top-$k$ index
selections were discrete; no additional prototype-specific loss or regional
labels were used. The bank indices do not identify validated MALT or SS
morphology, and a larger margin is not constrained to increase the final MALT score.

\paragraph{Coordinate-ordered groups and low-margin evidence.}
To form coordinate-ordered groups, let $\pi$ be the permutation obtained by
sorting the implemented coordinate key $10^6b_i+a_i$. The ordered hidden
sequence was partitioned using the code's chunk operation into $Q\leq16$
consecutive non-empty chunks $C_j$ (16 equal chunks for the 512 sampled patches).
Group tokens and the combined low-margin/regional vector were
\begin{subequations}\label{eq:regional}
\begin{align}
 \mathbf{r}_j &= \frac{1}{|C_j|}\sum_{l\in C_j}\mathbf{h}_{\pi(l)},
        \qquad j=1,\ldots,Q,\\
 \mathbf{e}_{\mathrm{reg}} &=
       \mathcal{G}_{\theta_R}(\{\mathbf{r}_j\}_{j=1}^{Q}),
       \qquad \mathbf{e}_{\mathrm{nr}}=\mathbf{e}_{-}+\mathbf{e}_{\mathrm{reg}}.
\end{align}
\end{subequations}
The regional operator includes its own value transformation from
Eq.~\eqref{eq:pool-value}. Coordinates determine group membership, but the
subsequent pooling is invariant to the order of the group tokens. The branch
has no explicit intergroup distances, adjacency or positional encoding.
Equal-sized groups of sampled instances were not anatomical regions and can
span different physical extents.

\paragraph{Similarity and margin statistics.}
Eight scalar statistics supplemented the evidence vectors. Define
$T_c=\operatorname{TopKIdx}(\{s_{i,c}\},k)$,
$\bar{s}_{c}^{\mathrm{top}}=k^{-1}\sum_{i\in T_c}s_{i,c}$ and
$\bar{m}=N^{-1}\sum_i m_i$. Then
\begin{subequations}\label{eq:statistics}
\begin{align}
 s_m &= \left[\frac{1}{N-1}\sum_{i=1}^{N}(m_i-\bar m)^2\right]^{1/2},\\
 \rho_0 &= \frac{1}{N}\sum_{i=1}^{N}\mathbb{I}(m_i>0),\qquad
 \rho_{\bar m}=\frac{1}{N}\sum_{i=1}^{N}\mathbb{I}(m_i>\bar m),\\
 \mathbf{u} &= \bigl[\bar{s}_{0}^{\mathrm{top}},\bar{s}_{1}^{\mathrm{top}},
       \bar m,s_m,\max_i m_i,\min_i m_i,\rho_0,\rho_{\bar m}\bigr]^{\mathsf T}
       \in\mathbb{R}^{8}.
\end{align}
\end{subequations}
$\mathbb{I}$ is the indicator function. The top-similarity sets $T_c$ are computed
separately for the two banks and need not coincide with the margin-based sets
$I_{+}$ and $I_{-}$.

\paragraph{Learned gating, weighted concatenation and classification.}
The five evidence vectors were ordered as
$(\mathbf{e}_1,\ldots,\mathbf{e}_5)=
(\mathbf{e}_{\mathrm{att}},\mathbf{e}_{\mu},\mathbf{e}_{\sigma},
 \mathbf{e}_{+},\mathbf{e}_{\mathrm{nr}})$.
A two-layer gate received the unweighted vectors and statistics:
\begin{subequations}\label{eq:fusion}
\begin{align}
 \mathbf{q} &= [\mathbf{e}_1;\mathbf{e}_2;\mathbf{e}_3;\mathbf{e}_4;
                   \mathbf{e}_5;\mathbf{u}]\in\mathbb{R}^{5D+8},\\
 \boldsymbol{\gamma} &= \operatorname{softmax}\!\left(
    \mathbf{W}_{g2}\operatorname{Dr}_{0.3}\!\left[
       \operatorname{GELU}(\mathbf{W}_{g1}\mathbf{q}+\mathbf{b}_{g1})\right]
       +\mathbf{b}_{g2}\right)\in\mathbb{R}^{5},\label{eq:gate}\\
 \widetilde{\mathbf{q}} &= [\gamma_1\mathbf{e}_1;\gamma_2\mathbf{e}_2;
       \gamma_3\mathbf{e}_3;\gamma_4\mathbf{e}_4;\gamma_5\mathbf{e}_5;
       \mathbf{u}],\label{eq:weighted-concat}\\
 \mathbf{z} &= \operatorname{Dr}_{0.3}\!\left[
       \operatorname{GELU}\!\left(\operatorname{LN}
       (\mathbf{W}_{f}\widetilde{\mathbf{q}}+\mathbf{b}_{f})\right)\right],\\
 a &= \mathbf{w}_{o}^{\mathsf T}\operatorname{Dr}_{0.3}\!\left[
       \operatorname{ReLU}\!\left(\operatorname{LN}
       (\mathbf{W}_{c}\mathbf{z}+\mathbf{b}_{c})\right)\right]+b_o,
       \qquad p=\operatorname{sigmoid}(a).
\end{align}
\end{subequations}
The gate dimensions were $\mathbf{W}_{g1}\in\mathbb{R}^{D\times(5D+8)}$
and $\mathbf{W}_{g2}\in\mathbb{R}^{5\times D}$;
$\mathbf{W}_{f}\in\mathbb{R}^{D\times(5D+8)}$ and
$\mathbf{W}_{c}\in\mathbb{R}^{D\times D}$.
Each $\gamma_j$ is non-negative and $\sum_j\gamma_j=1$.
Equation~\eqref{eq:weighted-concat} preserves five weighted feature blocks;
the eight statistics enter fusion without multiplication by a branch weight.
The gate supplies sample-dependent feature scaling, rather than a distribution
of causal contributions. Branch norms, the learned fusion transformation and
the ungated statistics also affect the classifier. The output $p$ is an
uncalibrated MALT model score in $[0,1]$.

\paragraph{Training objective.}
For a minibatch $\mathcal{B}$ of slides, the weighted binary cross-entropy was
\begin{equation}\label{eq:loss}
 \mathcal{L}=-\frac{1}{|\mathcal{B}|}\sum_{b\in\mathcal{B}}
 \left[w_{+}y_b\log p_b+(1-y_b)\log(1-p_b)\right],
 \qquad w_{+}=\frac{77}{67}.
\end{equation}
The implementation evaluated this objective directly from logits for numerical
stability. The positive-class weight was the ratio of the 77 SS to 67 MALT
training slides. Class weighting changes the optimum of binary cross-entropy;
the sigmoid output was not assumed to equal an unweighted posterior probability.
No calibration model was fitted to either evaluation set. Optimization and the
archived checkpoint are described in the main Methods.
